\section{Introducci\'on}
\label{sec:introduccion}

M\'exico tiene registradas 836 presas grandes y 4{,}330 presas peque\~nas, adem\'as de un estimado de 8{,}000 bordos que ni siquiera est\'an registrados \cite{infobae2020presas}. El 76\% del agua del pa\'is se destina a uso agr\'icola, buena parte de ella almacenada o derivada por esta infraestructura \cite{conagua2018estadisticas}. La Comisi\'on Nacional del Agua (CONAGUA) es responsable de vigilar la seguridad de esa red a trav\'es de su Sistema de Seguridad de Presas, pero la propia dependencia ha reconocido p\'ublicamente que desconoce el estado f\'isico y operativo de cerca de 1{,}000 presas por falta de presupuesto y personal para inspeccionarlas, incumpliendo con ello la recomendaci\'on de la Comisi\'on Internacional de Grandes Presas (ICOLD) de inspeccionar la red registrada al menos cada cinco a\~nos \cite{infobae2020presas, eluniversal2021presas}. De las presas s\'i evaluadas, 95 est\'an clasificadas actualmente como de alto riesgo \cite{infobae2020presas}. CONAGUA reconoce adem\'as que el riesgo de falla es ``permanente'', dado que la variabilidad hidrol\'ogica anual se ha vuelto m\'as abrupta por el cambio clim\'atico, adem\'as de fallas mec\'anicas, vandalismo y ausencia de mantenimiento \cite{infobae2020presas}.

Esta brecha no es un problema abstracto: buena parte de las presas peque\~nas y bordos sostienen el riego de comunidades agr\'icolas enteras, y muchas de las presas grandes del pa\'is llevan m\'as de 50 a\~nos en operaci\'on sin haber recibido mantenimiento mayor \cite{tecscience2023infraestructura}. La falla de una presa --por sismo, por una tormenta que excede la capacidad de dise\~no del vertedor, o simplemente por deterioro acumulado-- pone en riesgo directo a las poblaciones aguas abajo. El problema, en el fondo, es de asignaci\'on de un recurso escaso: si CONAGUA no tiene presupuesto para inspeccionar toda la red, la pregunta relevante no es \emph{si} se puede inspeccionar todo, sino \emph{qu\'e} se inspecciona primero.

Ese es exactamente el tipo de pregunta que un sistema de \emph{triage} basado en datos abiertos puede ayudar a responder, sin sustituir en ning\'un momento el criterio de un ingeniero especialista ni la inspecci\'on f\'isica en campo. La literatura en ingenier\'ia s\'ismica ya ofrece metodolog\'ias de evaluaci\'on r\'apida (\emph{screening}) para priorizar portafolios grandes de presas \cite{hariri2018prioritization, natural2023multiscale}, pero est\'an dise\~nadas y demostradas sobre inventarios que ya cuentan con datos estructurales b\'asicos (altura, geometr\'ia, a\~no de construcci\'on) capturados por el propio operador de la presa. Ese no es el caso de M\'exico: los atributos estructurales de la mayor parte del inventario, y en particular de las \(\sim\)1{,}000 presas sin supervisi\'on reciente que motivan este trabajo, no est\'an disponibles p\'ublicamente --viven \'unicamente en el buscador interactivo del Sistema de Seguridad de Presas, sin opci\'on de descarga masiva.

Este trabajo construye y eval\'ua, con ese vac\'io de datos como punto de partida expl\'icito, un \'indice de prioridad de revisi\'on basado \'unicamente en fuentes de datos abiertas y verificables: peligro s\'ismico, lluvia extrema hist\'orica, y poblaci\'on cercana. Se aplica sobre el subconjunto p\'ublico disponible --las 210 presas que CONAGUA monitorea activamente-- como validaci\'on metodol\'ogica, mientras se gestiona por separado, a trav\'es de una solicitud de transparencia gubernamental, el acceso al inventario completo y a los atributos estructurales necesarios para convertir este \'indice de exposici\'on en un screening de riesgo estructural pleno.

La Secci\'on \ref{sec:marco_teorico} revisa la literatura de evaluaci\'on de riesgo en portafolios de presas y ubica el vac\'io de datos mexicano frente a ella. La Secci\'on \ref{sec:metodologia} describe las fuentes de datos y la construcci\'on del \'indice. La Secci\'on \ref{sec:resultados} reporta su aplicaci\'on sobre las 210 presas monitoreadas. La Secci\'on \ref{sec:discusion} interpreta esos resultados y sus limitaciones. La Secci\'on \ref{sec:conclusiones} cierra con el alcance logrado y el trabajo pendiente.
